\documentclass[12pt,a4paper]{report}

\usepackage[utf8]{inputenc}
\usepackage[T1]{fontenc}
\usepackage[english]{babel}

\usepackage[top=2.5cm, bottom=2.5cm, left=3cm, right=2.5cm]{geometry}

\usepackage{microtype}

\usepackage{amsmath}

\usepackage{graphicx}
\usepackage{subfiles}

\usepackage{booktabs}
\usepackage{caption}
\usepackage{tabularx}

\usepackage{float}

\usepackage{enumitem}

\usepackage[table,dvipsnames]{xcolor}
\definecolor{codebg}{RGB}{248,248,248}
\definecolor{codeframe}{RGB}{200,200,200}
\definecolor{keywordcolor}{RGB}{0,0,200}
\definecolor{stringcolor}{RGB}{163,21,21}
\definecolor{commentcolor}{RGB}{0,128,0}

\usepackage{listings}

\lstdefinestyle{pythonStyle}{
    language=Python,
    backgroundcolor=\color{codebg},
    frame=single,
    rulecolor=\color{codeframe},
    basicstyle=\small\ttfamily,
    keywordstyle=\color{keywordcolor}\bfseries,
    stringstyle=\color{stringcolor},
    commentstyle=\color{commentcolor}\itshape,
    numberstyle=\tiny\color{gray},
    numbers=left,
    numbersep=8pt,
    stepnumber=1,
    showstringspaces=false,
    breaklines=true,
    breakatwhitespace=false,
    tabsize=4,
    captionpos=b,
}

\lstdefinestyle{jsonStyle}{
    language={},
    backgroundcolor=\color{codebg},
    frame=single,
    rulecolor=\color{codeframe},
    basicstyle=\small\ttfamily,
    stringstyle=\color{stringcolor},
    commentstyle=\color{commentcolor}\itshape,
    numberstyle=\tiny\color{gray},
    numbers=left,
    numbersep=8pt,
    stepnumber=1,
    showstringspaces=false,
    breaklines=true,
    breakatwhitespace=false,
    tabsize=2,
    captionpos=b,
    morestring=[b]",
    morecomment=[l]{//},
}

\lstset{style=pythonStyle}

\usepackage[
    colorlinks=true,
    linkcolor=NavyBlue,
    citecolor=NavyBlue,
    urlcolor=NavyBlue,
    pdftitle={Legal Contract Analysis: NLP Processing Pipeline},
    pdfauthor={HCMUT --- Semester 252},
    pageanchor=false,
]{hyperref}

\title{%
    \vspace{-1cm}
    \textbf{Legal Contract Analysis:}\\[6pt]
    \textbf{NLP Processing Pipeline}\\[12pt]
    \large Technical Report
}
\author{HCMUT\,---\,Semester 252 (2025\,--\,2026)}
\date{\today}

\begin{document}

\maketitle
\thispagestyle{empty}
\pagenumbering{roman}

\vspace{1cm}
\begin{abstract}
    This report documents a multi-stage NLP pipeline for analysing English legal
    contracts.  The system proceeds through three assignments.  Assignment~1
    implements a syntactic processing layer comprising clause splitting (spaCy
    sentence segmentation augmented with rule-based coordinating-conjunction and
    subordinating-conjunction detection), noun-phrase (NP) chunking with IOB
    tagging, and full Universal Dependencies parsing.  Assignment~2 adds a semantic
    extraction layer: a custom Named Entity Recognition (NER) model fine-tuned on
    \texttt{nlpaueb/legal-bert-base-uncased} with a rule-based fallback, a
    dependency-tree-based Semantic Role Labelling (SRL) module, and a dual-model
    intent classifier (TF-IDF\,+\,Logistic Regression and fine-tuned BERT).
    Assignment~3 integrates all upstream outputs into an end-to-end
    Retrieval-Augmented Generation (RAG) chatbot backed by ChromaDB, the
    \texttt{all-MiniLM-L6-v2} embedding model, Google Gemini, and a Streamlit web
    interface.
\end{abstract}

\newpage
\tableofcontents
\clearpage
\pagenumbering{arabic}
\hypersetup{pageanchor=true}

\subfile{assignment_1}

\newpage
\subfile{assignment_2}

\newpage
\subfile{assignment_3}

\newpage
\chapter{Discussion and Conclusion}
\label{chap:discussion}

\section{Summary of Achievements}
\label{sec:achievements}

The three assignments collectively deliver a complete, end-to-end NLP pipeline
for legal contract understanding.

\paragraph{Assignment 1 — Syntactic Pipeline.}
A full syntactic preprocessing layer was built using spaCy's
\texttt{en\_core\_web\_sm} model.  The clause splitter correctly handles
conditional sentences (preserving logical dependencies) and coordinates splits
(requiring both fragments to be grammatically complete).  IOB-tagged NP chunks
and full Universal Dependencies parse trees are produced for every clause,
providing the structural foundation for all downstream components.

\paragraph{Assignment 2 — Semantic Extraction.}
A three-layer semantic extraction stack was implemented.  The custom NER model
combines Legal-BERT fine-tuning with a rule-based fallback to identify six
domain-specific entity types (PARTY, MONEY, DATE, RATE, PENALTY, LAW).  The
dependency-tree-based SRL module extracts five semantic roles (Agent, Theme,
Recipient, Time, Condition) without requiring a computationally expensive dedicated
SRL model.  The dual-model intent classifier provides both an interpretable
TF-IDF baseline and a higher-accuracy BERT model across four intent classes.

\paragraph{Assignment 3 — RAG Chatbot.}
An interactive Streamlit application integrates all upstream outputs.  ChromaDB
with cosine similarity provides fast, semantically grounded clause retrieval.  The
LangChain LCEL chain orchestrates embedding, retrieval, and generation via Google
Gemini (\texttt{gemini-1.5-flash}), with an anti-hallucination system prompt that
mandates clause citation and explicit refusal when information is absent.

\section{Error Analysis}
\label{sec:error-analysis}

Despite solid overall performance, several systematic error patterns were observed
during qualitative evaluation.

\subsection*{Clause Splitting}
Long parenthetical enumerations (e.g.\ \textit{``Responsibilities include: (a)
filing reports, (b) attending meetings, (c) maintaining records''}) are sometimes
split at the enumeration markers rather than treated as a single obligation list.
Additionally, nested coordinating structures where an inner \texttt{CC} lacks a
clear local subject can produce over-splitting, creating non-sentential fragments
that should have been pruned but pass the two-token length guard.

\subsection*{NP Chunking}
Prepositional attachment ambiguity affects chunks for expressions such as
\textit{``payment of \$5,000 per month''}; depending on spaCy's parse, the PP
\textit{of \$5,000 per month} may be partially incorporated into the NP span,
inflating the chunk boundary.  Appositive NPs (e.g.\ \textit{``the Employer,
ABC Corp.''}) are treated as two separate chunks rather than one co-referential
span.

\subsection*{Dependency Parsing}
Passive-voice legal constructions (\textit{``shall be deemed to constitute a
material breach''}) and modal stacks (\textit{``shall have the right to terminate''})
sometimes produce unusual ROOT assignments.  In the latter example, spaCy may
select \textit{have} rather than \textit{terminate} as ROOT, misrepresenting the
logical predicate.

\subsection*{Named Entity Recognition}
The rule-based fallback misses multi-word party names that do not match known
patterns (e.g.\ \textit{``ABC International Co., Ltd.''} with trailing legal
suffixes).  Legal-BERT occasionally confuses RATE and MONEY labels for monetary
amounts that appear without explicit percentage signs (e.g.\ \textit{``at a rate
of three dollars per day''}).

\subsection*{Semantic Role Labelling}
The dependency-tree approach extracts roles only from the ROOT predicate, so
sentences with multiple verbs (e.g.\ \textit{``The Employer shall have the
obligation to pay and report''}) yield roles for only one predicate.  Stacked modal
constructions (\textit{``shall have the obligation to pay''}) cause the embedded
verb \textit{pay} to be missed as a secondary predicate entirely.

\subsection*{Intent Classification}
The word \textit{shall} in rights clauses (\textit{``Party A shall have the right
to inspect the premises''}) is misclassified as Obligation by both models because
\textit{shall} is the dominant Obligation indicator in the training data.  BERT
exhibits higher confidence on these errors, which is more dangerous than the
correctly low-confidence predictions from TF-IDF on the same clauses.

\section{Future Work}
\label{sec:future-work}

\subsection*{Clause Splitting}
Training a neural sentence-pair model (e.g.\ fine-tuned BERT) to classify whether
two consecutive clauses should be merged would address the enumeration-splitting
problem.  Enumerated list items (items starting with \textit{(a)}, \textit{(b)},
\textit{(i)}, etc.) should be detected by a pre-processing rule and treated as
atomic units before dependency-based splitting is applied.

\subsection*{Named Entity Recognition}
Annotating a larger domain-specific dataset with the six-entity schema would
substantially improve generalisation.  Exploring publicly available legal NER
benchmarks --- such as LexNER~\cite{luz2021lextreme} or CUAD~\cite{hendrycks2021cuad}
--- could provide additional fine-tuning signal.  Incorporation of a gazetteer for
known company-name suffixes (\textit{Ltd., Corp., S.A., GmbH}) would improve
recall for multi-token PARTY entities.

\subsection*{Semantic Role Labelling}
Integrating a proper BERT-based SRL model (e.g.\ AllenNLP's structured-prediction
SRL or a PropBank-fine-tuned sequence labeller) would handle multi-predicate
sentences and complex nested argument structures that the current dependency-tree
approach cannot resolve.

\subsection*{Intent Classification}
Adding a fifth class --- \textbf{Definition} --- for clauses that define terms
(e.g.\ \textit{``\,`Business Day' means any day other than a Saturday\ldots''})
would reduce label noise in the current four-class setting.  Experimenting with
legal domain-specific pre-training --- such as ContractBERT or LegalBERT variants
trained on larger contract corpora --- may further improve accuracy, particularly
on boundary cases between Right and Obligation.

\subsection*{RAG Chatbot}
Several enhancements would improve retrieval quality and answer fidelity: (i)
adding a cross-encoder re-ranker~\cite{nogueira2019passage} on top of the initial
top-$k$ retrieval to improve precision; (ii) extending the citation mechanism to
include precise character offsets rather than clause indices; (iii) supporting
multi-document contracts with document-level metadata filters; and (iv) evaluating
the chatbot's faithfulness and relevance using automated RAG evaluation frameworks
such as RAGAS~\cite{es2024ragas}.

\newpage
\begin{thebibliography}{9}

\bibitem{chalkidis2020legal}
I.~Chalkidis, M.~Fergadiotis, P.~Malakasiotis, N.~Aletras, and I.~Androutsopoulos,
\textit{LEGAL-BERT: The Muppets straight out of Law School},
in \textit{Findings of EMNLP 2020}, pp.\ 2898--2904, 2020.

\bibitem{hendrycks2021cuad}
D.~Hendrycks, C.~Burns, A.~Chen, and S.~Ball,
\textit{CUAD: An Expert-Annotated NLP Dataset for Legal Contract Review},
\textit{arXiv:2103.06268}, 2021.

\bibitem{luz2021lextreme}
P.~de~Araujo~Luz~de~Araujo, C.~Niklaus, and M.~Stürmer,
\textit{LEXTREME: A Multi-Lingual and Multi-Task Benchmark for the Legal Domain},
\textit{arXiv:2301.13126}, 2021.

\bibitem{nogueira2019passage}
R.~Nogueira and K.~Cho,
\textit{Passage Re-ranking with BERT},
\textit{arXiv:1901.04085}, 2019.

\bibitem{es2024ragas}
S.~Es, J.~James, L.~Espinosa-Anke, and S.~Schockaert,
\textit{RAGAS: Automated Evaluation of Retrieval Augmented Generation},
in \textit{EACL 2024}, 2024.

\end{thebibliography}

\end{document}
